\documentclass[CJK,13pt]{beamer}
\usepackage{CJKutf8}
\usepackage{beamerthemesplit}
\usetheme{Malmoe}
\useoutertheme[footline=authortitle]{miniframes}
\usepackage{amsmath}
\usepackage{amssymb}
\usepackage{graphicx}
\usepackage{eufrak}
\usepackage{color}
\usepackage{slashed}
\usepackage{simplewick}
\usepackage{tikz}
\usepackage{tcolorbox}
\usepackage{ulem}
\graphicspath{{../figures/}}
%%figures
\def\lfig#1#2{\includegraphics[width=#1 in]{#2}}
\def\tfig#1#2{\includegraphics[height=#1 in]{#2}}
\def\addfig#1#2{\begin{center}\includegraphics[width=#1 in]{#2}\end{center}}
\def\question{\includegraphics[width=0.3in]{why.jpg}\,}
\def\answer{\includegraphics[width=0.3in]{baozou_haha.png}\,}
\def\wulian{\includegraphics[width=0.18in]{emoji_wulian.jpg}}
\def\bigwulian{\includegraphics[width=0.35in]{emoji_wulian.jpg}}
\def\bye{\includegraphics[width=0.18in]{emoji_bye.jpg}}
\def\bigbye{\includegraphics[width=0.35in]{emoji_bye.jpg}}
\def\huaixiao{\includegraphics[width=0.18in]{emoji_huaixiao.jpg}}
\def\bighuaixiao{\includegraphics[width=0.35in]{emoji_huaixiao.jpg}}
\def\jianxiao{\includegraphics[width=0.18in]{emoji_jianxiao.jpg}}
\def\bigjianxiao{\includegraphics[width=0.35in]{emoji_jianxiao.jpg}}
\def\haoqi{\includegraphics[width=0.18in]{emoji_haoqi.jpg}}
%% colors
\def\blacktext#1{{\color{black}#1}}
\def\bluetext#1{{\color{blue}#1}}
\def\redtext#1{{\color{red}#1}}
\def\darkbluetext#1{{\color[rgb]{0,0.2,0.6}#1}}
\def\skybluetext#1{{\color[rgb]{0.2,0.7,1.}#1}}
\def\cyantext#1{{\color[rgb]{0.,0.5,0.5}#1}}
\def\greentext#1{{\color[rgb]{0,0.7,0.1}#1}}
\def\darkgray{\color[rgb]{0.2,0.2,0.2}}
\def\lightgray{\color[rgb]{0.6,0.6,0.6}}
\def\gray{\color[rgb]{0.4,0.4,0.4}}
\def\blue{\color{blue}}
\def\red{\color{red}}
\def\orange{\color[rgb]{1.,0.8,0.}}
\def\green{\color{green}}
\def\darkgreen{\color[rgb]{0,0.4,0.1}}
\def\darkblue{\color[rgb]{0,0.2,0.6}}
\def\skyblue{\color[rgb]{0.2,0.7,1.}}
%%control
\def\diag{\mathrm{diag}\,}
\def\heaviside{\,\mathrm{h}\,}
\def\bral{\left(\begin{array}{l}}
\def\brar{\end{array}\right)}
\def\brall{\left(\begin{array}{ll}}
\def\brarr{\end{array}\right)}
\def\bralll{\left(\begin{array}{lll}}
\def\brarrr{\end{array}\right)}
\def\branchl{\left\{\begin{array}{l}}
\def\branchr{\end{array}\right.}
\def\branchll{\left\{\begin{array}{ll}}
\def\branchrr{\end{array}\right.}
\def\branchlll{\left\{\begin{array}{lll}}
\def\branchrrr{\end{array}\right.}
\def\sfgamma#1{\,\Gamma\left( #1 \right)\,}
\def\be{\begin{equation}}
\def\ee{\nonumber\end{equation}}
\def\bea{\begin{eqnarray}}
\def\eea{\nonumber\end{eqnarray}}
\def\bch{\begin{CJK}{UTF8}{gbsn}}
\def\ech{\end{CJK}}
\def\bitem{\begin{itemize}}
\def\eitem{\end{itemize}}
\def\bcenter{\begin{center}}
\def\ecenter{\end{center}}
\def\bex{\begin{minipage}{0.2\textwidth}\includegraphics[width=0.6in]{jugelizi.png}\end{minipage}\begin{minipage}{0.76\textwidth}}
\def\eex{\end{minipage}}
\def\chtitle#1{\frametitle{\bch#1\ech}}
\def\bmat#1{\left(\begin{array}{#1}}
\def\emat{\end{array}\right)}
\def\bcase#1{\left\{\begin{array}{#1}}
\def\ecase{\end{array}\right.}
\def\bmini#1{\begin{minipage}{#1\textwidth}}
\def\emini{\end{minipage}}
\def\tbox#1{\begin{tcolorbox}#1\end{tcolorbox}}
\def\pfrac#1#2#3{\left(\frac{\partial #1}{\partial #2}\right)_{#3}}
\def\res#1#2{\,\mathrm{res}\,#1\left(#2\right)\,}
\def\newt#1#2{\left(\begin{array}{c}#1\\ #2\end{array}\right)}
\def\reof#1{\,\mathrm{Re}{\left(#1\right)\,}}
\def\imof#1{\,\mathrm{Im}{\left(#1\right)\,}}
\def\Arg#1{\,\mathrm{Arg}\,#1\,}
%%symbols
\def\intfull{\int_{-\infty}^\infty}
\def\inthalf{\int_0^\infty}
\def\bropt{\,(\ \ \ )}
\def\sech{\mathrm{sech}\,}
\def\csch{\mathrm{csch}\,}
\def\asinh{\mathrm{asinh}\,}
\def\acosh{\mathrm{acosh}\,}
\def\sone{$\star$}
\def\stwo{$\star\star$}
\def\sthree{$\star\star\star$}
\def\sfour{$\star\star\star\star$}
\def\sfive{$\star\star\star\star\star$}
\def\rint{{\int_\leftrightarrow}}
\def\roint{{\oint_\leftrightarrow}}
\def\stdHf{{\textit{\r H}_f}}
\def\deltaH{{\Delta \textit{\r H}}}
\def\ii{{\dot{\imath}}}
\def\skipline{{\vskip0.1in}}
\def\skiplines{{\vskip0.2in}}
\def\lagr{{\mathcal{L}}}
\def\hamil{{\mathcal{H}}}
\def\vecT{{\mathbf{T}}}
\def\vecN{{\mathbf{N}}}
\def\vecB{{\mathbf{B}}}
\def\vecv{{\mathbf{v}}}
\def\vecr{{\mathbf{r}}}
\def\vecf{{\mathbf{f}}}
\def\vecg{{\mathbf{g}}}
\def\vecupsilon{{\mathbf{\upsilon}}}
\def\vecu{{\mathbf{u}}}
\def\vecj{{\mathbf{j}}}
\def\vecx{{\mathbf{x}}}
\def\vecy{{\mathbf{y}}}
\def\vecz{{\mathbf{z}}}
\def\veck{{\mathbf{k}}}
\def\vecp{{\mathbf{p}}}
\def\vecn{{\mathbf{n}}}
\def\vecA{{\mathbf{A}}}
\def\vecP{{\mathbf{P}}}
\def\vecsigma{{\mathbf{\sigma}}}
\def\hatJn{{\hat{J_\vecn}}}
\def\hatJx{{\hat{J_x}}}
\def\hatJy{{\hat{J_y}}}
\def\hatJz{{\hat{J_z}}}
\def\hatj#1{\hat{J_{#1}}}
\def\hatphi{{\hat{\phi}}}
\def\hatq{{\hat{q}}}
\def\hatpi{{\hat{\pi}}}
\def\vel{\upsilon}
\def\Dint{{\mathcal{D}}}
\def\adag{{\hat{a}^\dagger}}
\def\bdag{{\hat{b}^\dagger}}
\def\cdag{{\hat{c}^\dagger}}
\def\ddag{{\hat{d}^\dagger}}
\def\hata{{\hat{a}}}
\def\hatb{{\hat{b}}}
\def\hatc{{\hat{c}}}
\def\hatd{{\hat{d}}}
\def\hatD{{\,\hat{D}}}
\def\hatN{{\hat{N}}}
\def\hatH{{\hat{H}}}
\def\hatp{{\hat{p}}}
\def\Fup{{F^{\mu\nu}}}
\def\Fdown{{F_{\mu\nu}}}
\def\newl{\nonumber \\}
\def\vece{\mathrm{e}}
\def\calM{{\mathcal{M}}}
\def\calT{{\mathcal{T}}}
\def\calR{{\mathcal{R}}}
\def\barpsi{\bar{\psi}}
\def\baru{\bar{u}}
\def\barv{\bar{\upsilon}}
\def\qeq{\stackrel{?}{=}}
\def\ftf{\stackrel{\mathcal{FT}}{\Longrightarrow}}
\def\ftb{\stackrel{\mathcal{FT}}{\Longleftarrow}}
\def\ftfb{\stackrel{\mathcal{FT}}{\Longleftrightarrow}}
\def\ltf{\stackrel{\mathcal{LT}}{\Longrightarrow}}
\def\ltb{\stackrel{\mathcal{LT}}{\Longleftarrow}}
\def\ltfb{\stackrel{\mathcal{LT}}{\Longleftrightarrow}}
\def\torder#1{\mathcal{T}\left(#1\right)}
\def\rorder#1{\mathcal{R}\left(#1\right)}
\def\contr#1#2{\contraction{}{#1}{}{#2}#1#2}
\def\trof#1{\mathrm{Tr}\left(#1\right)}
\def\trace{\mathrm{Tr}}
\def\comm#1{\ \ \ \left(\mathrm{used}\ #1\right)}
\def\tcomm#1{\ \ \ (\text{#1})}
\def\slp{\slashed{p}}
\def\slk{\slashed{k}}
\def\calp{{\mathfrak{p}}}
\def\veccalp{\mathbf{\mathfrak{p}}}
\def\Tthree{T_{\tiny \textcircled{3}}}
\def\pthree{p_{\tiny \textcircled{3}}}
\def\dbar{{\,\mathchar'26\mkern-12mu d}}
\def\erf{\mathrm{erf}}
\def\const{\mathrm{const.}}
\def\pheat{\pfrac p{\ln T}V}
\def\vheat{\pfrac V{\ln T}p}
%%units
\def\fdeg{{^\circ \mathrm{F}}}
\def\cdeg{^\circ \mathrm{C}}
\def\atm{\,\mathrm{atm}}
\def\angstrom{\,\text{\AA}}
\def\SIL{\,\mathrm{L}}
\def\SIkm{\,\mathrm{km}}
\def\SIyr{\,\mathrm{yr}}
\def\SIGyr{\,\mathrm{Gyr}}
\def\SIV{\,\mathrm{V}}
\def\SImV{\,\mathrm{mV}}
\def\SIeV{\,\mathrm{eV}}
\def\SIkeV{\,\mathrm{keV}}
\def\SIMeV{\,\mathrm{MeV}}
\def\SIGeV{\,\mathrm{GeV}}
\def\SIcal{\,\mathrm{cal}}
\def\SIkcal{\,\mathrm{kcal}}
\def\SImol{\,\mathrm{mol}}
\def\SIN{\,\mathrm{N}}
\def\SIHz{\,\mathrm{Hz}}
\def\SIm{\,\mathrm{m}}
\def\SIcm{\,\mathrm{cm}}
\def\SIfm{\,\mathrm{fm}}
\def\SImm{\,\mathrm{mm}}
\def\SInm{\,\mathrm{nm}}
\def\SImum{\,\mathrm{\mu m}}
\def\SIJ{\,\mathrm{J}}
\def\SIW{\,\mathrm{W}}
\def\SIkJ{\,\mathrm{kJ}}
\def\SIs{\,\mathrm{s}}
\def\SIkg{\,\mathrm{kg}}
\def\SIg{\,\mathrm{g}}
\def\SIK{\,\mathrm{K}}
\def\SImmHg{\,\mathrm{mmHg}}
\def\SIPa{\,\mathrm{Pa}}
%page
\def\secpage#1#2{\begin{frame}\bcenter{\bf \Huge #1} \ecenter \skipline \tbox{\bcenter #2 \ecenter}\end{frame}}
\def\append#1#2{\secpage{附录#1}{#2}}
\def\thinka#1{\begin{frame}\frametitle{思考题}\bcenter\lfig{0.4}{think0.jpg}  \skipline #1 \ecenter\end{frame}}
\def\thinkb#1{\begin{frame}\frametitle{思考题}\bcenter\lfig{0.5}{think1.jpg}  \skipline #1 \ecenter\end{frame}}
\def\thinkc#1{\begin{frame}\frametitle{思考题}\bcenter\lfig{0.5}{think2.jpg}  \skipline #1 \ecenter\end{frame}}
\def\thinkd#1{\begin{frame}\frametitle{思考题}\bcenter\lfig{0.5}{think3.jpg}  \skipline #1 \ecenter\end{frame}}
\def\thinke#1{\begin{frame}\frametitle{思考题}\bcenter\lfig{0.5}{think4.jpg}  \skipline #1 \ecenter\end{frame}}
\def\thinkf#1{\begin{frame}\frametitle{思考题}\bcenter\lfig{0.5}{think5.jpg}  \skipline #1 \ecenter\end{frame}}
\def\schw{ds^2 = \left(1-\frac{2GM}{r}\right)dt^2 - \left(1-\frac{2GM}{r}\right)^{-1} dr^2 - r^2\left(d\theta^2 + \sin^2\theta d\phi^2\right)}

\input{titlepage.tex}
  \date{}
  \begin{document}
  \bch
\tpage{2}{Embedded Curves}


\begin{frame}
  \frametitle{欧式空间中的正则曲线}
  设$n$维欧氏空间有条曲线 $\vecx(t) = \left(x_1(t), x_2(t), \ldots, x_n(t)\right)$。

  \skipline

  它的切向量可以定义为 $\frac{d\vecx}{dt} = \left(\frac{dx_1}{dt}, \frac{dx_2}{dt}, \ldots, \frac{dx_n}{dt}\right)$，这也是一个向量函数。

  \skipline
  
  本讲中我们将一直假设曲线的各个分量函数至少三次可导，且切向量处处非零： $\left\vert\frac{d\vecx}{dt}\right\vert\ne 0$。这样的曲线称之为“{\blue 正则曲线}”。

\end{frame}



\begin{frame}
  \frametitle{曲线的弧长参数和单位切向量$\vecT$}
  如果参数从某个允许的确定值 $a$ (例如，如果允许的话，可取 $a=0$) 变化到 $t$， 对应的一段弧长为
  $$ s = \int_a^t \left\vert\frac{d\vecx}{d\tau}\right\vert d\tau. $$
  我们可以选取 $s$ 而不是 $t$ 作为参量来描述曲线。
  
  {\scriptsize 注1：正则性保证了 $t$ 到 $s$ 的映射是单调映射。}

  \skipline

  {\scriptsize 注2：一般我们假定曲线的参数形式默认给了曲线一个方向，即 $t$ 的正向给出了曲线的方向。选定方向的曲线的弧长参数 $s$ 几乎就是确定的了，至多相差一个常数。我们后面讨论的 Frenet 标架等，都依赖于曲线方向已经给定这个假定。否则曲线的弧长参数和 Frenet 标架会有个符号不确定性。}

  \skipline
  
  选择了弧长参数之后，显然地， $\vecT \equiv \frac{d\vecx}{ds}$ 是{\blue 单位切向量}。  
\end{frame}


\begin{frame}
  \frametitle{曲率(curvature)}

  曲线的一小段总是能用圆弧来近似。圆弧半径越大，曲线的一小段就看起来越直。反之，圆弧半径越小，曲线的一小段就看起来很弯曲。所以曲线在一点上的{\blue “曲率” $\kappa$ 就定义为这点附近的近似圆弧的半径的倒数}。
  
  \addfig{2}{curvature.jpg}

\end{frame}



\begin{frame}
  \frametitle{曲线的曲率公式}

  假设有曲线 $\vecx(s)$， $s$ 是弧长参数。在曲线上的给定点附近， 单位切向量 $\frac{d\vecx}{ds}$ 的变化的大小应该等于它的方向变化的大小（用近似圆弧来计算就是 $\frac{ds}{R}$）。
  
  $$ \left\vert d \frac{d\vecx}{ds}\right\vert = \frac{ds}{R} $$


  \addfig{2}{curvature_explain.jpg}

\end{frame}

\begin{frame}
  \frametitle{曲线的曲率公式}
  由此得到曲率 $\frac{1}{R}$ 的值
  {\blue
  $$\kappa =  \left\vert\frac{d^2\vecx}{ds^2}\right\vert. $$}

  这个公式如此简洁美妙，完全得益于我们选取了弧长作为曲线的参数。

  \skipline
  
  如果我们把取绝对值的操作去掉，则得到的是一个“曲率向量”：
  $$\frac{d^2\vecx}{ds^2}. $$
  它的大小为曲率，方向从所讨论的点指向近似圆弧的圆心。  


\end{frame}


\begin{frame}
  \frametitle{密切平面, 主法向量，次法向量}
  假设曲线上某点的曲率向量 $\frac{d^2\vecx}{ds^2}$ 非零 (显然它和该点单位切向量 $\vecT = \frac{d\vecx}{ds}$ 垂直)， 
  我们称归一化的曲率向量，也就是：
  $$ \vecN \equiv \frac{1}{\kappa} \frac{d^2\vecx}{ds^2}$$
  为曲线在该点的{\blue 主法向量}(normal vector)。


  单位切向量 $\vecT$ 和主法向量 $\vecN$ 一起确定了近似圆弧所在的“{\blue 密切平面}”； 和密切平面垂直的 $n-2$ 个单位法向量则称为{\blue 次法向量} (binormal vectors)。
  
\end{frame}

\begin{frame}
  \frametitle{三维空间曲线的Frenet标架}


  从现在开始我们专门讨论嵌入在三维空间的曲线。这时单位切向量$\vecT$，主法向量$\vecN$，次法向量$\vecB$一起构成了一个局域的直角坐标架——它有个很炫酷的名字，叫“{\blue Frenet 标架}”。  

  \addfig{2.7}{FrenetFrame.png}
  

\end{frame}


\begin{frame}
  \frametitle{嵌在三维空间的曲线的Frenet标架总结}  
  切向量(tangent vector)
  $$\vecT(s) = \frac{d\vecx}{ds};$$
  主法向量(normal vector)
  $$\vecN(s) = \frac{1}{\kappa}\frac{d^2\vecx}{ds^2};\ \kappa \equiv \lvert\frac{d^2\vecx}{ds^2}\rvert;$$
  次法向量(binormal vector)
  $$\vecB(s) = \vecT(s)\times \vecN(s).$$

  {\scriptsize 注：垂直于$\vecT(s)$, $\vecN(s)$, $\vecB(s)$ 的平面分别称为法平面，切平面，和密切平面。}
\end{frame}



\begin{frame}
  \frametitle{嵌在三维空间的曲线的挠率(torsion)}
  一般来说当沿着曲线运动时，密切平面会发生“摆动”，即次法向量 $\vecB$ 会发生变化。我们先来证明，$\frac{d\vecB}{ds}$ 一定和 $\vecN$ 同向。

  把归一化条件 $\vecB^2 = 1$ 对 $s$ 求导，可得到 $\frac{d\vecB}{ds} \cdot \vecB = 0$，即 $\frac{d\vecB}{ds}$ 和 $\vecB$ 是垂直的。此外，
  $$\frac{d\vecB}{ds} = \frac{d(\vecT\times \vecN)}{ds} = \frac{d\vecT}{ds}\times \vecN + \vecT \times \frac{d\vecN}{ds}$$
  注意到 $\frac{d\vecT}{ds} = \kappa \vecN$，上式右边的第一项为零，即 $\frac{d\vecB}{ds}$ 和 $\vecT$ 也垂直。
  
  于是不妨设
  $$\frac{d\vecB}{ds} = -\tau \vecN.$$
  我们把 $\tau$ 称为{\blue 挠率}。注意和曲率不同，挠率是可正可负的。它的大小 $|\tau|= \lvert \frac{d\vecB}{ds}\rvert$ 刻画了密切平面的摆动快慢。
\end{frame}

\thinkf{如果一条曲线的挠率处处为零，这说明了曲线的什么几何特性？}


\begin{frame}
  \frametitle{嵌在三维空间的曲线的Frenet方程}
  我们来研究Frenet标架的变化率。首先，根据定义已经知道了：
  $$\frac{d\vecT}{ds} = \kappa \vecN$$
  和
  $$\frac{d\vecB}{ds} = -\tau \vecN.$$  
  剩下的非常简单：
  $$\frac{d\vecN}{ds} = \frac{d\vecB}{ds}\times \vecT + \vecB\times\frac{d\vecT}{ds} =-\tau\vecN\times\vecT +\kappa\vecB\times\vecN = \tau\vecB - \kappa \vecT.$$
  
\end{frame}


\begin{frame}
  \frametitle{把Frenet方程写成矩阵形式}
  \be
  \frac{d}{ds}
  \begin{pmatrix}
    \vecT \\
    \vecN \\
    \vecB
  \end{pmatrix}
  =
  \begin{pmatrix}
    0 & \kappa & 0 \\
    -\kappa & 0 & \tau \\
    0 & -\tau & 0
  \end{pmatrix}
  \begin{pmatrix}
    \vecT \\
    \vecN \\
    \vecB
  \end{pmatrix}
  \ee
\end{frame}


\begin{frame}
  \frametitle{思考题（本题给出了挠率的一种常见计算方法）}
  设 $\vecx(s)$ 是三维空间的一条曲线，$s$是弧长参数。试证明：以三个矢量 $\frac{d\vecx}{ds}$, $\frac{d^2\vecx}{ds^2}$, $\frac{d^3\vecx}{ds^3}$ 为列矢量的行列式的值等于 $\tau \kappa^2$。
\end{frame}


\begin{frame}
  \frametitle{如果曲线不是用弧长参数表示的怎么办}
  一般参数表示的曲线 $\vecx(t)$ 的Frenet标架、曲率、挠率的计算都比较复杂。但都可以用
  $$\frac{ds}{dt} = \lvert\frac{d\vecx}{dt}\rvert$$
  进行转化。下面我们直接给出计算结果：
\end{frame}


\begin{frame}
  \frametitle{一般参数曲线$\vecx(t)$的Frenet计算公式}
  \begin{eqnarray}
    \vecT(t) &=& \frac{\vecx'(t)}{\lvert\vecx'(t)\rvert} \newl
    \vecN(t) &=& \frac{\lvert\vecx'(t)\rvert^2\vecx''(t)-\left(\vecx'(t)\cdot\vecx''(t)\right)\vecx'(t)}{\lvert\vecx'(t)\rvert\,\lvert\vecx'(t)\times\vecx''(t)\rvert } \newl    
    \vecB(t) &=& \frac{\vecx'(t)\times \vecx''(t)}{\lvert\vecx'(t)\times\vecx''(t)\rvert} \newl
    \kappa(t) &=& \frac{\lvert\vecx'(t)\times\vecx''(t)\rvert}{\lvert\vecx'(t)\rvert^3} \newl
    \tau(t) &=& \frac{\vecx'(t)\cdot(\vecx''(t)\times\vecx'''(t))}{\lvert\vecx'(t)\times\vecx''(t)\rvert^2} \nonumber 
  \end{eqnarray}

  请参考 \url{http://zhiqihuang.top/gr/codes/frenet.py}
\end{frame}


\begin{frame}
  \frametitle{曲线的邻域展开}
  由于 $\vecT$, $\vecN$, $\vecB$ 包含了曲线的一次导数，二次导数，三次导数的信息，我们可以在一点附近展开曲线到三阶近似：

  \bea
  \vecx(s+ds) & \approx & \vecx(s) \newl
  && + \left(ds - \frac{\kappa^2}{6}ds^3\right)\vecT(s) \newl
  &&  + \left(\frac{\kappa}{2}ds^2 + \frac{1}{6}\frac{d\kappa}{ds}ds^3\right)\vecN(s) \newl
  && + \frac{\kappa\tau}{6}ds^3\vecB(s). 
  \eea  
\end{frame}


\begin{frame}
  \frametitle{例题：球面上的曲线}
  设 $\vecx(s)$ 是三维空间的球面上的一条曲率和挠率均非零的曲线，$s$ 是弧长参数。证明：
  $$\frac{1}{\kappa^2} +\left[\frac{1}{\tau}\frac{d}{ds}\left(\frac{1}{\kappa}\right)\right]^2$$
  是常数。

\end{frame}


\begin{frame}
  \frametitle{解法1}
  证明：不妨设球半径为单位长度，并取球心为原点。球面上的曲线满足方程：
  \begin{equation}
  \vecx^2 = 1.
  \end{equation}
  两边对 $s$ 求导得到
  \begin{equation}
  \vecx\cdot\vecT = 0.
  \end{equation}
  于是可以假设
  \begin{equation}
  \vecx = \alpha \vecN + \beta \vecB
  \end{equation}
  其中 $\alpha$, $\beta$ 是待定的（依赖于$s$的）参量，满足 $\alpha^2+\beta^2=1$。两边继续对 $s$ 求导并利用Frenet公式，
  \begin{equation}
    \vecT =-\alpha\kappa\vecT + \left(\frac{d\alpha}{ds}-\tau\beta\right)\vecN + \left(\frac{d\beta}{ds} + \alpha\tau\right)\vecB.
  \end{equation}
\end{frame}

\begin{frame}
  \frametitle{解法1}
  于是有
  \begin{eqnarray}
    \alpha &=& -\frac{1}{\kappa};  \\
    \beta &=& -\frac{1}{\tau}\frac{d}{ds}\left(\frac{1}{\kappa}\right).
  \end{eqnarray}
  于是利用 $\alpha^2+\beta^2=1$ 即得证。
\end{frame}


\begin{frame}
  \frametitle{解法2}
  证明：不妨设球半径为单位长度，并取球心为原点。球面上的曲线满足方程：
  \begin{equation}
  \vecx^2 = 1.
  \end{equation}
  把 $\vecx(s+ds)^2 = 1$ 近似展开到 $ds^3$量级，有
  \bea
  1 &=& 1 + (2ds -\frac{\kappa}{3}ds^3)\vecx(s)\cdot\vecT(s) \newl
  && +\left(\kappa ds^2 + \frac{1}{3}\frac{d\kappa}{ds}ds^3\right)\vecx(s) \cdot \vecN(s) \newl
  && + \frac{\kappa\tau}{3}ds^3 \, \vecx(s)\cdot\vecB(s) \newl
  && + ds^2.
  \eea  

\end{frame}

\begin{frame}
  \frametitle{解法2}
  由 $ds$一次项系数为零，得到
  $$\vecx(s)\cdot\vecT(s)=0.$$  
  由 $ds^2$系数为零，得到
  $$\vecx(s)\cdot\vecN(s) = -\frac{1}{\kappa}.$$
  再由 $ds^3$ 系数为零，得到
  $$\vecx(s)\cdot\vecB(s) = \frac{1}{\kappa^2\tau}\frac{d\kappa}{ds} = -\frac{1}{\tau}\frac{d}{ds}\left(\frac{1}{\kappa}\right).$$
  于是得证。
\end{frame}


\begin{frame}
  \frametitle{曲线论基本定理}
  根据我们对曲率和挠率的几何图像理解，下列定理是显然的：

  \skipline
  
  {\blue 三维欧氏空间内两条以弧长 $s$ 为参数的正则曲线。如果它们的曲率处处不为零，且对同一个 $s$，两条曲线的曲率，挠率都分别相等。则两条曲线可以通过一个刚体运动重叠起来。}

\end{frame}


\begin{frame}
  \frametitle{例题}
  证明曲线 $\vecx(t) = (\cosh t, \sinh t, t)$ 和曲线 $\vecy(u) = \left(\frac{e^{-u}}{\sqrt{2}}, \frac{e^u}{\sqrt{2}}, u+1\right)$ 在三维欧式空间内的一个刚体运动下是重合的。
\end{frame}



\begin{frame}
  \frametitle{解答}
  证明: 直接计算得到 $\vecx(t)$ 弧长参数，曲率和挠率为：
  $$ s_x(t) = \sqrt{2} \sinh t, \kappa_x(t) = \tau_x(t) = \frac{1}{2\cosh^2(t)} $$
  $\vecy(u)$ 的弧长参数，曲率和挠率为
  $$ s_y(u)  = \sqrt{2} \sinh u, \kappa_y(u) = \tau_y(u) = \frac{1}{2\cosh^2(u)}.$$
  故同样的 $t$ 和 $u$ 对应相同的弧长参数，曲率，挠率。证毕。
\end{frame}


\begin{frame}
  \frametitle{二维欧式空间中的曲线}
  嵌入在二维欧式空间里的曲线可以看成嵌在三维欧式空间中的无挠（挠率为零的）曲线。平面曲线有一些独特的性质，我们来专门讨论一下。

  设曲线 $\vecr(s) = (x(s), y(s))$ 以弧长 $s$ 为参数。则切向量为
  $$ \vecT(s) = \left(x'(s), y'(s)\right),$$
  逆时针方向旋转$\frac{\pi}{2}$，得到{\blue 相对法向量}为
  $$\vecN_r(s) = \left(-y'(s), x'(s)\right).$$
  注意三维情况定义 $\vecN(s)$ 时我们是按“往那边弯曲就指向哪边”的原则来定义的。所以 $\vecN_r$ 和 $\vecN$ 有可能相同，也可能差一个负号。
  
\end{frame}



\begin{frame}
  \frametitle{相对曲率}
  由相对法向量可以定义``{\blue 相对曲率}''：
  $$ \kappa_r \equiv \vecN_r \cdot \frac{d^2\vecx}{ds^2} $$
  显然它可能和三维情况定义的 $\kappa$ 相同，也可能差一个负号。

  容易看出来，相对曲率可以写成行列式的形式
  \be
  \kappa_r = \left\vert
  \begin{array}{ll}
    x'(s) & y'(s) \\
    x''(s) & y''(s)     
  \end{array}
  \right\vert
  \ee
  它可正可负，而之前定义的 $\kappa = |\kappa_r|$ 一定非负.
\end{frame}

\begin{frame}
  \frametitle{相对曲率}
  对一般的参数形式表示的平面曲线 $\left(x(t), y(t)\right)$，相对曲率可以写成
  \be
  \kappa_r = \frac{1}{\left[x'(t)^2+ y'(t)^2\right]^{3/2}}\left\vert
  \begin{array}{ll}
    x'(t) & y'(t) \\
    x''(t) & y''(t)     
  \end{array}
  \right\vert
  \ee
\end{frame}


\begin{frame}
  \frametitle{平面曲线的 Frenet 方程}
  平面曲线的 Frenet 方程也变得非常简单：

  \be
  \frac{d}{ds}
  \begin{pmatrix}
    \vecT \\
    \vecN_r
  \end{pmatrix}
  =
  \begin{pmatrix}
    0 & \kappa_r \\
    -\kappa_r & 0 
  \end{pmatrix}
  \begin{pmatrix}
    \vecT \\
    \vecN_r
  \end{pmatrix}
  \ee
\end{frame}


\ech
\end{document}
